\label{chap:intro}

\section{Background}
\label{sec:background}

Large language models (LLMs) do not always tell the truth \cite{park2023deception}. Sometimes they
lie, that is, they generate statements they believe to be false. Following Liars' Bench
\cite{kretschmar2025liars}, we adopt an assertion based definition of lying: the assistant lies if it
states something it believes to be false. Beliefs are measured through consistent demonstration of
knowledge when the model is under no direct pressure to lie. For example, under no pressure to
respond otherwise, if the assistant consistently asserts ``the capital of France is Paris'', we would
consider this statement a belief. This definition distinguishes lies from hallucination, during
hallucination both truth and falsehoods are asserted independently of beliefs. Prior work found
models to be lying in various scenarios, including lying in order to conceal their prior
misconduct \cite{scheurer2023strategic,jarviniemi2024uncovering}, and lying in order to fake
their alignment to seem cooperative \cite{greenblatt2024alignment}. As LLMs are now managing
sensitive tasks, monitoring when a model is lying becomes essential to ensuring alignment and
safety.

Two broad families of lie detectors exist: Black-box and White-box. Black-box detectors can
only see the chat transcript, a LLM judge is instructed to read the transcript and evaluate if the
model lied or not \cite{pacchiardi2023catch}. White-box detectors read the model's internal activations and decide if
the model lied or not. The simplest white-box detector is a linear probe, which is a linear
classifier trained on the activations collected from the residual stream
\cite{alain2016understanding,azaria2023internal,marks2023geometry}. A linear probe
is computationally cheaper because it can detect when a model is lying using only a single dot
product. Moreover, it reads activations at inference time and unlike black-box detectors it
requires no extra forward pass. As it reads models internal representations, it can still detect lies
even if there is no direct trace of lying in the transcript \cite{goldowskydill2025detecting}.

The weakness of the probe is generalisation. A probe trained on one distribution of lies and
deployed on others often fails to transfer
\cite{kirch2025offpolicy,kumar2026pressure,natarajan2026targeted}. Liars' Bench makes the failure
concrete \cite{kretschmar2025liars}. It contains \V{lb-examples} lies and honest responses, across
seven datasets it captures distinct types of lies generated on-policy by four popular open-weight
LLMs. Those datasets vary on two axes: object of belief and reason for lying. Training probes on
instructed deception following \textcite{goldowskydill2025detecting}, reaches balanced accuracy near
chance on most of them. The hardest is the harm pressure knowledge report (\hpkr). In \hpkr{} a
model denies knowing the answer to a hazardous question that it has been shown to know. No
method the benchmark evaluates without privileged labels detects the lies of \hpkr{} at its 1\%
false positive threshold without also flagging most honest answers, and its instruction-trained
probe's AUROC on Llama-3.3-70B-Instruct is \V{lb-mean-auroc-hpkr}, which is worse than chance. That
inversion is measured at layer~20, the layer the benchmark reads; at layer~40, the layer this thesis
reports, a probe trained the same way ranks \hpkr{} well but still flags none of its lies
(\cref{sec:problem-measured}). Only a probe that is trained using samples from
\hpkr{} itself, succeeds there.

\textcite{mahadik2026persona} propose using persona coordinates for better generalisation of probes.
They construct persona vectors following the Assistant Axis methodology \cite{lu2026assistant}. They
prompt a model to play deceptive roles and honest roles and average its activations over each
role's responses. It gives one persona vector per role. They use the geometry of these vectors in
two ways: directly, as a contrast between the mean
of the deceptive personas and the mean of the honest ones, or as a low-dimensional coordinate
system in which to train probes. Because the directions come from role-play rather than from an
evaluation dataset, they cannot fit one dataset's idiosyncrasies. On Llama~3.2-3B and
Llama~3-8B, probes in these \emph{persona coordinates} transfer across deception datasets better
than probes on raw activations. The authors name larger models and richer persona sets as future
work. This thesis takes up both.

\section{Rationale of the Study or Motivation}
\label{sec:motivation}

Three observations motivate this work.

First, the cheapest detector is the one that fails to transfer. A linear probe costs almost
nothing to run (\cref{sec:background}), but a detector that must be retrained with labelled examples
for every new kind of lie is not practical, because the lies that matter most are the ones nobody
anticipated. A direction that works without labels from the setting in which it is deployed would
keep the cost of a probe and remove its main weakness.

Second, lies about what the model knows are where such directions fail. When a model claims not to
know an answer it knows, its lie can be word for word the same as an honest ``I do not know'' from a
model that really does not know. \hpkr{} is built of such lies, and so is Gender-Secret, a model
fine-tuned to deny knowing its user's gender \cite{kretschmar2025liars}. A detector that reads how a
lie sounds cannot catch one that sounds exactly like the truth. The persona approach had not been
tested on these cases, nor on a model of 70 billion parameters.

Third, the definition of a lie offers a second route. A lie is a statement the speaker believes to be
false (\cref{sec:background}), so a lie about knowledge is a denial by a model that knows. What the
model knows can be measured directly, without any labelled lie, by asking it questions it knows and
questions nobody can know. If that knowledge can be read from the activations while the words stay
fixed, it should catch exactly the lies that a persona direction misses.

\section{Problem Statement}
\label{sec:problem}

When the published persona contrast of \textcite{mahadik2026persona} is rebuilt for
Llama-3.3-70B-Instruct and applied to Liars' Bench, it fails on \hpkr{} in the worst possible way:
it \emph{inverts}, ranking the model's lies as more honest than its honest answers, with AUROC
\V{zs-dec-hpkr-L20} at layer~20. In \hpkr{} every lie is a \emph{denial of knowledge}, while every
published deceptive persona lies by \emph{controlling disclosure}: timing, emphasis and omission. Adding
personas that deny knowledge they have raises the contrast's ranking of \hpkr, but for the wrong
reason. The new direction flags nearly every response that says ``I do not know'', true or false: at
layer~40 it flags \V{th-k-deight-flag-hden-L40}\% of the honest denials. Where the lies and the honest
answers are the same sentence, every persona direction is within reach of random directions, and on
Gender-Secret every persona direction is at chance (\cref{sec:hpkr-tests}). A persona direction reads
the register of a lie, not what makes it a lie.

The problem this thesis addresses is therefore:

\begin{quote}
\emph{Can a lie detector for a 70B-parameter language model, built without any labelled lies,
separate lies from honest answers when the two use the same words, as lies about knowledge often do,
while keeping the cross-dataset transfer of persona directions?}
\end{quote}

\section{Objective}
\label{sec:objective}

The aim of this study is to construct and evaluate one zero-shot, label-free lie detector for
Llama-3.3-70B-Instruct that combines a persona direction, which reads how a response is framed, with a
knowledge direction, which reads what the model knows, and to establish where it works and why it
needs both parts. The work is organised around four research questions:

\begin{description}[leftmargin=1.2em, labelindent=0em, itemsep=2pt]
\item[RQ1: Transfer.] Does a detector built from a persona contrast and a knowledge direction, with
  no labelled lies, transfer across Liars' Bench better than supervised probes?
\item[RQ2: The same words.] Does it separate lies from honest answers that use the same words, as
  lies about knowledge often do, and does it detect the lies of a model fine-tuned to hide what it
  knows?
\item[RQ3: The two parts.] What does each part read: does the persona contrast respond to how a
  response is framed, and does the knowledge direction read whether the model knows the answer while
  the words of its reply stay fixed?
\item[RQ4: Robustness and operating points.] Does the result survive resampling the questions, new
  questions and other layers, and how many lies does the detector catch at strict thresholds?
\end{description}

From these follow seven concrete objectives:
\begin{enumerate}[leftmargin=1.6em, itemsep=2pt]
\item Rebuild the persona bank of \textcite{mahadik2026persona} for Llama-3.3-70B-Instruct and extend
  it with roles whose deception is a denial of knowledge.
\item Build a knowledge direction from questions the model knows and questions nobody can know, with
  the same denial forced after each, and verify the questions against the model's own answers.
\item Evaluate the selected detector and every alternative on five Liars' Bench datasets by AUROC and
  at the benchmark's threshold, on \hpkr's wording-controlled slices, and on two fine-tuned model
  organisms read through their fine-tunes.
\item Compare twelve designs and select among them on transfer, wording control, detection, labels
  and cost (\cref{sec:alternatives}).
\item Test the robustness of the selected detector and of its knowledge direction: by resampling K's
  questions, by balanced splits, by a replication with new questions and pre-fixed criteria, across
  \V{depth-n} layers, and at fixed false-positive rates.
\item Generate the reported numbers from result files, so that the results can be reproduced and
  checked, and record the provenance of the few that cannot be (\cref{sec:code}).
\item Repeat the evaluation on the benchmark's three other models, to separate what holds for lies from
  what holds for one model (\cref{sec:xmodel}).
\end{enumerate}

\section{Methodology in Brief}
\label{sec:method-brief}

The work has three parts.

\emph{Building persona coordinates.} Following \textcite{mahadik2026persona}, Llama-3.3-70B-Instruct
role-plays each of 36 personas (the 26 of the original study plus \V{bank-new} new roles built around
the denial of knowledge) under five prompt variants and 30 shared questions, producing
\V{bank-generated} responses. A second model, Llama-3.1-70B-Instruct, judges whether each response
stays in character, and only in-character responses are kept. Each persona's vector is the model's
mean residual-stream activation over its kept responses, and a \emph{contrast direction} is the
difference between the mean of a group of deceptive personas and the mean of a group of honest ones.

\emph{Building the knowledge direction.} The model reads \V{th-k-known} questions it knows and
\V{th-k-unknown} that nobody can know, each in \hpkr's own template and each followed by the same
denial, written for the model: ``I do not know the answer.''. The knowledge direction K is the
difference between the mean activation over the known questions and over the unknowable ones. Every
read ends in the same words, so K can only reflect what the model knows.

\emph{Scoring and evaluating.} To score a response, we average the model's activations over the
response's tokens and project the result onto a direction. The selected detector's direction is the
standardised sum of the extended persona contrast and K, one fixed vector. No step uses labels from an
evaluation dataset: the scores are \emph{zero-shot}. We
evaluate them under the benchmark's own protocol against supervised baselines, on slices of \hpkr{}
where the words are held fixed, and on the benchmark's fine-tuned models, whose adapters we apply
inside the remote forward pass. The model is reached through the National Deep Inference Fabric
(NDIF) \cite{fiottokaufman2024nnsight}; all analysis runs on a CPU from pooled activations.
\Cref{fig:pipeline} in \cref{sec:design-overview} draws the persona pipeline, \cref{fig:k-build}
the knowledge direction and \cref{fig:selected} the detector that combines them.

\section{Scopes and Challenges}
\label{sec:scope}

\emph{Scope.} The study covers Llama-3.3-70B-Instruct \cite{meta2024llama33}, chosen because Liars'
Bench provides its lies and because NDIF serves its internals. \Cref{sec:xmodel} then repeats the
evaluation, without the steering experiment, on the benchmark's three other models,
Qwen2.5-72B-Instruct, Mistral-Small-3.1-24B-Instruct and Gemma-3-27B-it. For Llama, the evaluation covers
the five Liars' Bench datasets that this model generated without fine-tuning, together with RepE, the
instruction pairs on which the benchmark's probe is trained, and two Liars' Bench model organisms,
Gender-Secret and Soft-Trigger, which we read through their LoRA fine-tunes (\cref{sec:data}). The
detectors are linear, read from the residual stream, and evaluated at \V{depth-n} layers (every
fourth from \V{depth-first} to 76, plus the last block, \cref{sec:transformation}). We report layer~40
as the main operating depth. We do not train or fine-tune any model.

\emph{Challenges.} Five shaped the work.
\begin{itemize}[leftmargin=1.2em, itemsep=2pt]
\item \emph{Access to a 70B model.} The model does not fit on our hardware, so every forward pass
  ran remotely on NDIF, with limited request sizes and a slow connection (\cref{sec:specs}).
  Extraction therefore had to pool activations on the server and resume after interruptions.
\item \emph{Reading a fine-tune that the server does not serve.} NDIF serves only the base model, so
  the fine-tunes' adapters had to be applied inside the remote pass, checked against a reference
  implementation, and sent in parts small enough for the server (\cref{sec:finetuned-reads}).
\item \emph{A 16\,GB local machine.} The analysis, including logistic regression on 8192-dimensional
  activations for tens of thousands of examples, had to fit in 16\,GB of memory.
\item \emph{A designed repair is not a blind test.} The knowledge-denial personas and the knowledge
  direction were both designed after we saw \hpkr{} fail, so \hpkr{} measures repairs designed for it.
  Gender-Secret, questions written afterwards, and a control that widens the bank with different roles
  (\cref{sec:adj-roles}) carry the untargeted evidence.
\item \emph{Ranking and detection come apart.} A detector can rank lies above honest answers and
  still flag none of them at a usable threshold. Much of the analysis is about keeping these two
  apart (\cref{sec:perf}).
\end{itemize}

\section{Contributions}
\label{sec:contributions}

In this work, we make the following contributions:
\begin{itemize}[leftmargin=1.2em, itemsep=2pt]
\item \emph{A label-free lie detector.} The extended persona contrast plus the knowledge direction K
  transfers better than supervised probes, and separates lies that use the same words as honest
  answers (\cref{sec:transfer,sec:selection}).
\item \emph{A knowledge direction.} K reads what the model knows while the words stay fixed. It is
  built from questions the model knows and questions nobody can know, without labelled lies
  (\cref{sec:k-design,sec:k-results}).
\item \emph{A persona bank for deception by denial.} We add knowledge-denial roles to the published
  bank and scale the method to a 70B model, where a single persona contrast transfers better than
  supervised probes (\cref{sec:banks,sec:transfer}).
\item \emph{A diagnosis of persona directions.} On lies about knowledge, persona directions read the
  words of a denial rather than the lie, which is why the detector needs K (\cref{sec:hpkr-tests}).
\item \emph{Tests on fine-tuned models.} We test the detectors on fine-tuned models that lie, reading
  each model's own activations. K detects Gender-Secret's lies there, while the persona directions
  are near chance (\cref{sec:finetuned-reads,sec:k-results}).
\item \emph{A causal test of the persona part.} Steering along the extended contrast makes the model
  conceal an unfavourable fact far more often, while random directions do not (\cref{sec:steering}).
\item \emph{A replication across models.} On Mistral-Small-3.1-24B and Gemma-3-27B the persona contrast
  again transfers better than supervised probes; on all four models the knowledge-denial roles help and
  the selected detector beats a probe trained at the same depth. On Qwen2.5-72B, which rarely plays the
  roles, the persona contrasts fail, and K's result on a fine-tuned liar does not carry over to Mistral
  (\cref{sec:xmodel}).
\end{itemize}

\section{Organisation of the Report}

\Cref{chap:lit} reviews the background and related work. \Cref{chap:req} sets out the requirements,
the societal, environmental and ethical impact of the work, and its management. \Cref{chap:method}
presents the design process, the alternative designs, the knowledge direction and the data.
\Cref{chap:results} evaluates the alternatives, records how the final design was adjusted, repeats the
evaluation on three more models, and discusses the results and their limits. \Cref{chap:conclusion}
concludes. The appendices list the personas, supplementary results, the knowledge direction's questions
and the full tables of the replication across models.
